\section{System Design and Methodology}

\subsection{System Requirements and Overview}
The system is organized into two pipelines. The offline pipeline prepares legal evidence. It validates curated legal texts, splits them into hierarchy-aware chunks, enriches metadata, extracts references, builds the Qdrant index, and constructs the Neo4j legal graph. The online pipeline answers user queries. It retrieves seed chunks, expands them through graph relations and fallbacks, reranks and budgets context, constructs an answer, validates citations, and returns either a grounded answer or insufficient-context response.

\thesisfigure{overall_system_architecture.pdf}{Overall System Architecture.}{fig:overall-system-architecture}

The design requirements follow directly from the legal QA setting: preserve statutory coordinates, retrieve both semantic and exact-coordinate evidence, use graph relations without allowing uncontrolled expansion, validate all final citations against retrieved context, and refuse questions whose legal basis is outside the indexed corpus.

\subsection{Legal Corpus and Hierarchy-Aware Processing}
The corpus is intentionally scoped to six legal document groups. This scope supports reproducible evaluation and bounded system claims. Table~\ref{tab:corpus-summary-final} is based on \path{artifacts/index/current.json}.

\begin{table}[H]
  \centering
  \caption{Scoped Legal Corpus Summary}
  \label{tab:corpus-summary-final}
  \small
  \begin{tabularx}{\linewidth}{@{}L{0.30\linewidth}L{0.30\linewidth}rr@{}}
    \toprule
    Document ID & Role in Corpus & Rank & Chunks \\
    \midrule
    \path{45-2019-qh14} & Labor Code 2019 & 1 & 697 \\
    \path{92-2015-qh13-labor-only} & Civil Procedure Code labor subset & 1 & 159 \\
    \path{nghi-dinh-135-2020-nd-cp} & Retirement age decree & 2 & 41 \\
    \path{nghi-dinh-145-2020-nd-cp} & Labor Code implementation decree & 2 & 520 \\
    \makecell[l]{\texttt{thong-tu-09-2020-}\\\texttt{tt-bldtbxh}} & Minor worker circular & 3 & 73 \\
    \makecell[l]{\texttt{thong-tu-10-2020-}\\\texttt{tt-bldtbxh}} & Labor contract circular & 3 & 66 \\
    \midrule
    Total & Six indexed document groups & & 1,556 \\
    \bottomrule
  \end{tabularx}
\end{table}

The rank column encodes the legal hierarchy used by the system. Rank 1 contains primary codes, rank 2 contains decrees, and rank 3 contains circulars. This hierarchy does not decide legal correctness by itself, but it helps context assembly avoid treating lower-level guidance as a substitute for primary statutory authority.

The system uses curated text files rather than raw OCR during the final pipeline. The cleaning and validation stage normalizes text, checks required fields, and preserves official document boundaries. The chunking stage is hierarchy-aware: it detects articles, clauses, points, and appendices instead of cutting by arbitrary character windows. Each chunk stores both text for retrieval and text for citation. This separation matters because search benefits from parent headings and expanded context, while citations need stable human-readable labels.

The key metadata fields include identifiers, document labels, normative rank, article and clause coordinates, citation text, retrieval text, and taxonomy labels for topic, actor, and issue type. These fields are carried through Qdrant payloads and used later by validation.

\subsection{Hybrid Retrieval and Legal Knowledge Graph}
For a query $q$, Qdrant returns an initial seed set:
\[
S_k = \operatorname{HybridSearch}(q,k).
\]
Hybrid search combines dense Vietnamese SBERT vectors with sparse lexical matching. Dense vectors support conceptual similarity. Sparse signals support exact coordinates and legal vocabulary. The final index manifest records \texttt{keepitreal/vietnamese-sbert}, 768-dimensional dense vectors, a sparse vector channel, and 1,556 indexed records.

The graph is designed for provision-level expansion. It represents evidence chunks, legal documents, articles, clauses, points, appendices, topics, actors, and issue types. Edges represent structural hierarchy, source links, cross-references, taxonomy membership, and normative hierarchy. The graph is deterministic: it is built from processed artifacts and rules rather than from open-ended LLM extraction.

At runtime, seed chunks are mapped to Neo4j nodes through their shared chunk identifiers. The graph expander then adds related contexts:
\[
C_{\mathrm{graph}} = \operatorname{Expand}_{G}(S_k, d, E),
\]
where $d$ is the configured expansion depth and $E$ is the set of allowed edge types. Direct coordinate fallback also retrieves missing explicit article references when possible. The final context set is:
\[
C_{\mathrm{final}} =
\operatorname{Rerank}(S_k \cup C_{\mathrm{graph}} \cup C_{\mathrm{fallback}}).
\]

\thesisfigure{graph_augmented_retrieval_flow.pdf}{Graph-Augmented Retrieval Flow.}{fig:graph-augmented-retrieval-flow}

\subsection{Answer Construction and Citation Validation}
The answer layer is deliberately conservative. It receives only budgeted retrieved contexts. It can use provider-backed generation when configured, but the official final evaluation uses a deterministic extractive provider. The validator checks citation containment: every legal basis cited in the answer must be available in the retrieved context, and article numbers mentioned in the answer must be supported by retrieved metadata. If a generated answer fails validation, the system can return an extractive fallback or insufficient-context response.

\subsection{Evaluation Design}
The evaluation separates retrieval, answer quality, citation validation, and refusal behavior. Table~\ref{tab:metrics-summary-final} summarizes the metrics used in Section~5.

\begin{table}[H]
  \centering
  \caption{Evaluation Metrics Summary}
  \label{tab:metrics-summary-final}
  \small
  \begin{tabularx}{\linewidth}{@{}L{0.25\linewidth}L{0.18\linewidth}Y@{}}
    \toprule
    Metric & Scope & Meaning \\
    \midrule
    Recall@10 & Retrieval & Fraction of required citations recovered in the top-10 retrieved contexts. \\
    Required Citation Coverage & Retrieval & Citation-level coverage over in-corpus queries. \\
    Forbidden Citation Violation Rate & Retrieval & Share of in-corpus queries where a benchmark-defined distracting citation appears. \\
    Retrieval pass rate & Retrieval & Query-level pass requiring required citations and no forbidden violation. \\
    Answer pass rate & Answer & Share of answers passing deterministic answer-presence and rule checks. \\
    Citation validation pass rate & Answer & Share of answers whose citations remain inside retrieved context. \\
    Quality validation pass rate & Answer & Share of answers passing deterministic quality checks. \\
    Adjusted end-to-end pass rate & Full QA & In-corpus QA pass plus out-of-corpus refusal pass over all 100 queries. \\
    Out-of-corpus QA pass rate & Safety & Share of unsupported queries refused or marked insufficient. \\
    \bottomrule
  \end{tabularx}
\end{table}
